\section{Introduction}\label{sec:intro}

The medical field is one of the domains where impressive amounts of data are
collected, ranging from general patient information (age, body mass index, blood
type) to detailed clinical records such as treatment plans, prescribed
medications, surgical interventions, immunizations, observations,
hospitalizations and allergies. Properly processed, this data can help
clinicians identify \emph{patterns that repeat across patients}, reveal groups
that require similar care, and surface the treatments that worked best for each
group. In practice, however, much of this data either remains on paper or sits
untouched in hospital databases. Its potential is enormous but, as Yadav
\emph{et al.}~\cite{Yadav2018MiningEHR} observe, largely wasted.

Building a system that closes this gap raises four coupled challenges. The
\textbf{first} is obtaining and preprocessing the heterogeneous source data. The
\textbf{second} is grouping patients with a method that can combine
\emph{numerical} attributes (e.g., age, number of active conditions) with
\emph{categorical} ones (e.g., gender, dominant care-plan category, comorbidity
flags) in a single model. The \textbf{third} is understanding the
\emph{temporal} structure of care through sequential pattern mining. The
\textbf{fourth} is making the resulting analytics interpretable, and the
database queryable, for users without a technical background.

Because real EHR data is hard to obtain and is protected by privacy
regulations, we use the open-source Synthea generator~\cite{Walonoski2018Synthea}
to produce a realistic synthetic cohort of \NPATIENTS{} patients, each with a
deep simulated medical history spanning organizations, providers, encounters,
conditions, allergies, observations, immunizations, care plans, procedures and
medications. The depth of the inter-table relationships makes this dataset a
suitable instrument for all four challenges.

To address them we introduce \textbf{HospitAI}, a hospital analytics platform
organized as a single layered pipeline. Patients are clustered with
K-Prototypes~\cite{Huang1998KPrototypes}, which extends K-Means with a combined
distance that uses squared Euclidean distance for numerical features and simple
matching for categorical ones. Care trajectories are discovered with
PrefixSpan~\cite{Pei2004PrefixSpan}, a projection-based sequential pattern miner.
Crucially, we apply PrefixSpan both globally and \emph{per cluster}, so that
group-characteristic trajectories can be distinguished from population-wide
ones. Finally, a large language model (LLM) layer turns numerical outputs into
natural language and exposes a guarded text-to-SQL interface. The use of
unsupervised learning to discover patient clusters in EHRs is itself a promising
direction in recent literature~\cite{Wang2020UnsupervisedEHR}; our emphasis is
on \emph{integration} and \emph{interpretability} rather than on a new learning
algorithm.

\smallskip
\noindent\textbf{Contributions.} This paper makes the following contributions:
\begin{itemize}
  \item An end-to-end, interpretation-first pipeline that reduces raw multi-table
        EHR data to one mixed-type feature vector per patient and feeds it
        through clustering, sequence mining and an LLM layer
        (Sections~\ref{sec:overview}--\ref{sec:llm}).
  \item A \emph{dual-scope} application of PrefixSpan that mines clinical
        sequences both globally and independently within each patient cluster,
        together with a frequency-corrected ranking based on geometric-mean lift
        and a clinical type-diversity bonus (Section~\ref{sec:sequences}).
  \item An LLM layer that operates at every analytical step --- cluster
        summarization, sequence interpretation and a validate-then-execute
        text-to-SQL agent with role-based scoping
        (Section~\ref{sec:llm}).
  \item A qualitative and quantitative evaluation on a \NPATIENTS{}-patient
        synthetic cohort showing interpretable clusters and cluster-specific
        trajectories (Section~\ref{sec:results}).
\end{itemize}

The remainder of the paper is organized as follows.
Section~\ref{sec:related} reviews related work.
Section~\ref{sec:overview} describes the system architecture.
Section~\ref{sec:ingestion}--\ref{sec:llm} detail the methods.
Section~\ref{sec:results} reports results, Section~\ref{sec:discussion}
discusses limitations, and Section~\ref{sec:conclusion} concludes.
